\documentclass{beamer}

\usepackage{amsmath}
\usepackage{booktabs}
\usepackage{graphicx}
\usepackage{makecell}

\title{Spatio-temporal EEG effect recovery}
\subtitle{Hard-threshold simulation and SEM BGAM}
\author{Jiahui Fan}
\date{\today}

\graphicspath{
    {../results/figures/n20_hard_sem_slides/}
    {../results/figures/}
    {../pred_vs_raw_compare/}
}

\newcommand{\missingfigure}[1]{%
    \fbox{\parbox{0.82\linewidth}{\centering
    Figure not found. Run \texttt{#1} first.}}}

\begin{document}

\frame{\titlepage}

% ----------------------------------------------------------------
\begin{frame}{Question}
\small
Can a spatio-temporal BGAM recover a known EEG effect over the full
channel $\times$ time grid?

\medskip
Current simulation:
\begin{itemize}
\item $P = 20$ participants;
\item $R = 20$ trials per participant and condition;
\item $61$ electrodes;
\item true signal window: $160$--$342$ ms;
\item seed electrode: CZ;
\item condition 1 = noise, condition 2 = noise + signal.
\end{itemize}
\end{frame}

% ----------------------------------------------------------------
\begin{frame}{Signal model}
\small
For electrode $c$ and time $t$:
\[
y_{c,t}^{(1)} = \epsilon_{c,t}^{(1)}
\]
\[
y_{c,t}^{(2)}
= A \, g(t) \, w(c) + \epsilon_{c,t}^{(2)}
\]

\medskip
\begin{itemize}
\item $A = 1.5\,\mu\mathrm{V}$;
\item $g(t)$ is the ERP template with true window $160$--$342$ ms;
\item $w(c)$ is the CZ-centred spatial weight;
\item noise combines temporal EEG-like noise and spatially correlated sensor
noise.
\end{itemize}
\end{frame}

% ----------------------------------------------------------------
\begin{frame}{Spatial ground truth: active electrodes}
\small
The current simulation uses a hard-threshold spatial Gaussian.

\[
d(c, \mathrm{CZ})
= \left\|p_c - p_{\mathrm{CZ}}\right\|_2
\]

\[
w(c) =
\begin{cases}
\exp\!\left(-\dfrac{d(c,\mathrm{CZ})^2}
{2\sigma_{\mathrm{space}}^2}\right),
& d(c,\mathrm{CZ}) \leq \sigma_{\mathrm{space}},\\[0.7em]
0,
& d(c,\mathrm{CZ}) > \sigma_{\mathrm{space}}.
\end{cases}
\]

\[
\text{electrode } c \text{ is active iff }
d(c,\mathrm{CZ}) \leq \sigma_{\mathrm{space}}.
\]

\medskip
This is the definition used in \texttt{code/00\_generate\_data.R} and in
the ground-truth active-electrode list used for F1, FP, and FN metrics.
\end{frame}

% ----------------------------------------------------------------
\begin{frame}[fragile]{BGAM model}
\scriptsize
\begin{verbatim}
bgam_fit <- brm(
    formula = eeg_mean | se(eeg_sem) ~
        condition +
        (1 + condition | participant) +
        t2(xproj, yproj, time,
           by = condition,
           bs = c("tp", "tp"),
           k = c(k_xy, k_t),
           d = c(2, 1),
           full = full),
    data = sim_data_model,
    family = gaussian(),
    chains = 4,
    cores = 4,
    iter = 3000,
    warmup = 1000,
    threads = threading(threads = 4),
    backend = "cmdstanr",
    control = list(adapt_delta = 0.99, max_treedepth = 13)
)
\end{verbatim}
\end{frame}

% ----------------------------------------------------------------
\begin{frame}{Single-simulation workflow}
\small
\begin{enumerate}
\item Load one generated dataset and its ground truth.
\item Fit the SEM BGAM.
\item Predict both conditions on the full channel $\times$ time grid.
\item Compute posterior probability and odds for the condition difference.
\item Compare methods using connected time-channel detection masks.
\item Save one CSV row per method.
\end{enumerate}

\medskip
Implemented in \texttt{code/09\_run\_montecarlo\_sem.R}.
\end{frame}

% ----------------------------------------------------------------
\begin{frame}{Simulated EEG data ($P=20$, $R=20$)}
\begin{center}
\IfFileExists{../results/figures/n20_hard_sem_slides/sim_data_overview_n20_hard_sem.pdf}
{\includegraphics[width=\linewidth]
{../results/figures/n20_hard_sem_slides/sim_data_overview_n20_hard_sem.pdf}}
{\missingfigure{code/05\_slide\_simulated\_eeg\_n20.R}}
\end{center}
\end{frame}

% ----------------------------------------------------------------
\begin{frame}{How the model fits the data ($P=20$, $k=(20,20)$)}
\begin{center}
\IfFileExists{../results/figures/n20_hard_sem_slides/pred_vs_raw_bgam_fit_v5_k20_sem.pdf}
{\includegraphics[height=0.85\textheight,keepaspectratio]
{../results/figures/n20_hard_sem_slides/pred_vs_raw_bgam_fit_v5_k20_sem.pdf}}
{\missingfigure{code/06\_slide\_bgam\_fit\_n20.R}}
\end{center}
\end{frame}

% ----------------------------------------------------------------
\begin{frame}{Posterior probability per electrode}
\begin{center}
\IfFileExists{../results/figures/n20_hard_sem_slides/posterior_p_sig_facets_n20_hard_sem.pdf}
{\includegraphics[width=0.92\linewidth]
{../results/figures/n20_hard_sem_slides/posterior_p_sig_facets_n20_hard_sem.pdf}}
{\missingfigure{code/07\_slide\_posterior\_probability\_n20.R}}
\end{center}
\end{frame}

% ----------------------------------------------------------------
\begin{frame}{Topomaps over the signal window}
\begin{center}
\IfFileExists{../results/figures/n20_hard_sem_slides/topomap_n20_hard_sem_three_times.pdf}
{\includegraphics[height=0.86\textheight,keepaspectratio]
{../results/figures/n20_hard_sem_slides/topomap_n20_hard_sem_three_times.pdf}}
{\missingfigure{code/08\_slide\_topomap\_n20.R}}
\end{center}
\end{frame}

% ----------------------------------------------------------------
\begin{frame}{Method comparison}
\small
Each method returns a time $\times$ channel detection mask.

\medskip
\begin{enumerate}
\item Raw pointwise $t$-test.
\item Benjamini--Hochberg FDR.
\item Holm correction.
\item MNE spatio-temporal cluster permutation.
\item MNE TFCE.
\item BGAM posterior odds on the full channel $\times$ time grid.
\end{enumerate}

\medskip
For each method, small isolated blobs are removed and the largest connected
blob gives onset and offset.
\end{frame}

% ----------------------------------------------------------------
\begin{frame}{Key output: pilot model comparison}
\scriptsize
\textbf{Pilot model} here means the early smaller setting with
$P=5$ participants and $R=20$ trials per condition. It was used to compare
model specifications before the current $P=20$ hard-threshold SEM run.

\medskip
\resizebox{\linewidth}{!}{%
\begin{tabular}{ccccrrrrrrr}
\toprule
$k_{xy}$ & $k_t$ & \texttt{full} & SE &
\makecell{Onset\\bias (s)} & \makecell{Offset\\bias (s)} &
\makecell{Peak\\rec.\,(\%)} & \makecell{Mean\\rec.\,(\%)} &
\makecell{Spatial\\F1} & LOOic & \makecell{Wall-\\clock} \\
\midrule
12 & 12 & F & SD  & $-0.070$ & $+0.068$ & 46 & 55.4  & $0.714$ & 57462.5 & 5\,h \\
20 & 20 & F & SD  & $-0.070$ & $+0.068$ & 46 & 54.9  & $0.714$ & 57465.5 & 14\,h \\
20 & 20 & F$^{\star}$ & SD & $-0.080$ & $+0.078$ & 40 & 75.5 & $0.758$ & 57462.4 & 28\,h \\
\textbf{20} & \textbf{20} & \textbf{F} & \textbf{SEM} &
$\mathbf{+0.010}$ & $\mathbf{-0.012}$ & \textbf{89} & \textbf{106.1} &
$\mathbf{0.962}$ & -- & 67\,h \\
10 & 30 & T & SD  & $-0.080$ & $+0.068$ & 36 & 54.3 & $0.714$ & 57469.4 & 48\,h \\
20 & 30 & T & SD  & $-0.130$ & $+0.108$ & 40 & 63.2 & $0.658$ & 57468.7 & 112\,h \\
30 & 20 & T & SD  & $-0.150$ & $+0.138$ & 58 & 57.7 & $0.633$ & 57490.4 & 69\,h \\
30 & 30 & T & SD  & $-0.150$ & $+0.138$ & 36 & 57.7 & $0.625$ & 57489.9 & 110\,h \\
\bottomrule
\end{tabular}}

\smallskip
\textsuperscript{$\star$} Best SD fit in the original pilot sweep.
\end{frame}

% ----------------------------------------------------------------
\begin{frame}{Key output: methods comparison ($P=5$ pilot, 30 MC)}
\scriptsize
Thirty independent simulations under the $P=5$, $R=20$ pilot setting.
Onset/offset are in seconds; SD is across simulations.

\medskip
\textbf{Onset estimates} (truth $=0.160$\,s)

\resizebox{\linewidth}{!}{%
\begin{tabular}{lrrrrc}
\toprule
Method & Bias & MAE & RMSE & SD & $F1_{\text{ST}}$ (mean $\pm$ SD) \\
\midrule
Raw $p$       & $+0.010$ & $\mathbf{0.010}$ & $0.012$ & $0.008$ & $0.462 \pm 0.014$ \\
\textbf{BGAM\_2D} & $-0.011$ & $0.011$ & $\mathbf{0.012}$ & $\mathbf{0.005}$ & $\mathbf{0.799 \pm 0.007}$ \\
TFCE          & $+0.013$ & $0.014$ & $0.017$ & $0.011$ & $0.539 \pm 0.105$ \\
Cluster       & $+0.016$ & $0.016$ & $0.017$ & $0.007$ & $0.514 \pm 0.020$ \\
FDR (BH)      & $+0.063$ & $0.063$ & $0.065$ & $0.016$ & n/a \\
Holm          & $+0.093$ & $0.093$ & $0.094$ & $0.012$ & n/a \\
\bottomrule
\end{tabular}}

\medskip
\textbf{Offset estimates} (truth $=0.342$\,s)

\resizebox{0.82\linewidth}{!}{%
\begin{tabular}{lrrrr}
\toprule
Method & Bias & MAE & RMSE & SD \\
\midrule
\textbf{BGAM\_2D} & $+0.012$ & $\mathbf{0.012}$ & $\mathbf{0.014}$ & $\mathbf{0.006}$ \\
Raw $p$       & $-0.006$ & $0.013$ & $0.015$ & $0.014$ \\
TFCE          & $-0.015$ & $0.015$ & $0.018$ & $0.011$ \\
Cluster       & $-0.018$ & $0.019$ & $0.020$ & $0.009$ \\
FDR (BH)      & $-0.068$ & $0.068$ & $0.070$ & $0.017$ \\
Holm          & $-0.082$ & $0.082$ & $0.082$ & $0.010$ \\
\bottomrule
\end{tabular}}
\end{frame}

% ----------------------------------------------------------------
\begin{frame}{Key output: FP/FN diagnosis ($\alpha=0.10$)}
\scriptsize
Onset/offset metrics summarise only the first and last detected time points.
The full time $\times$ channel mask shows the false-positive and
false-negative trade-off.

\medskip
\resizebox{\linewidth}{!}{%
\begin{tabular}{lrrrrr}
\toprule
Method & FP rate & FN rate & $F1_{\text{ST}}$ &
\makecell{Mean FP\\cells} & \makecell{FP-only\\clusters} \\
\midrule
Raw $p$  & $0.102$ & $0.530$ & $0.462$ & $268.9$ & $141.4$ \\
BH       & $0.001$ & $0.943$ & $0.106$ & $3.0$   & $2.9$ \\
Holm     & $0.000$ & $0.996$ & $0.008$ & $0.2$   & $0.2$ \\
Cluster  & $0.003$ & $0.648$ & $0.514$ & $8.0$   & $0.0$ \\
TFCE     & $0.007$ & $0.608$ & $0.539$ & $18.4$  & $1.1$ \\
\bottomrule
\end{tabular}}

\medskip
Raw $p$ detects early boundaries well but creates many scattered FP cells.
BH and Holm control FP but miss most true cells. Cluster and TFCE use
spatio-temporal structure, giving a better FP/FN trade-off.
\end{frame}

% ----------------------------------------------------------------
\begin{frame}{Final aggregation}
\small
After all Monte Carlo runs finish:

\[
\texttt{code/04\_aggregate\_results.R}
\]

\medskip
The final table contains:
\begin{itemize}
\item onset and offset bias;
\item MAE, RMSE, and SD;
\item spatio-temporal F1;
\item TP, FP, and FN diagnostics.
\end{itemize}

\medskip
Main output:
\[
\texttt{results/comparisons/<MC\_AGG\_PREFIX>\_paper\_table.csv}
\]
\end{frame}

% ================================================================
\appendix

\begin{frame}{Appendix: FP/FN definitions}
\small
Each cell is one time-by-channel location $(t,c)$.

\[
M_{t,c} =
\begin{cases}
1, & t \in [0.160, 0.342]\text{ s and } c \text{ is active},\\
0, & \text{otherwise}.
\end{cases}
\]

\[
D_{t,c} =
\begin{cases}
1, & \text{method detects cell }(t,c),\\
0, & \text{otherwise}.
\end{cases}
\]

\[
TP=\sum D_{t,c}M_{t,c},\quad
FP=\sum D_{t,c}(1-M_{t,c}),\quad
FN=\sum (1-D_{t,c})M_{t,c}.
\]
\end{frame}

% ----------------------------------------------------------------
\begin{frame}{Appendix: prior predictive check}
\small
\texttt{code/10\_prior\_predictive\_check.R} runs the current SEM BGAM formula
with \texttt{sample\_prior = "only"}.

\medskip
A weak normal prior is added only because prior predictive sampling requires
a proper fixed-effect prior.
\end{frame}

% ----------------------------------------------------------------
\begin{frame}{Appendix: files not part of the main story}
\small
The main story is the hard-threshold + SEM Monte Carlo pipeline.

\medskip
Details moved out of the main slide sequence:
\begin{itemize}
\item old pilot-only diagnostics;
\item old SD-based model comparisons;
\item LOO notes;
\item threshold-sweep experiments;
\item exploratory smoothing-bias notes.
\end{itemize}
\end{frame}

\end{document}
